In this section, I will build intuition for the role that each component plays by speaking generally about narratives using semantics used by humans in english. I will use it to demonstrate how the narrative can be constructed in a structured form and how that meaning can be used to construct more meaning. Then I will demonstrate how the meaning of narratives can be reasoned about to identify if it is internally consistent or if it has incorrect or missing meaning. Then I will discuss how the components outlined in this document work toward the same goal for software systems. 

Consider the following narrative and participants in the narrative. In this case, the world is defined with the librarian and shelf. The user interacts with the world through the librarian and narratives unfold. 

\subsubsection*{Librarian}
\begin{greyquote}
    The librarian presents the user with the option to add a book or get a book from the library.\\

    If the user selects the option to add a book, the librarian accepts a book by asking the user for the name and content. Then the librarian reads the first letter of the books name to place the book on the relevant shelf. Then the user is once again asked which action they want to take.\\

    If the user selects the option to get a book, the librarian asks the user for the book name. The librarian reads the first letter of the book name and accesses the book from the relevant shelf. After the librarian provides the user the book, the user is once again asked which action they want to take.
\end{greyquote}

\subsubsection*{Book}
\begin{greyquote}
    A book has a name and it has contents. 
\end{greyquote}

\subsubsection*{Shelf}
\begin{greyquote}
    A shelf has slots for books. The slots are labelled using letters.
\end{greyquote}






